% Procedencia íntegra: CONTRAANGULO_ELECTRON_BARBERO_CATALAN.md, §§2--10.
% La gestión, el historial y los localizadores se conservan en gestion/.
\section{Composition et restitution des lectures angulaire, électronique et constitutive}
\label{md:compos:seccion}

\subsection{Coefficient d’action et construction du contre-angle}
\label{md:compos:contraangulo}

L’abréviation fonctionnelle $B(\alpha)$ réunit une évaluation en $\alpha$ ; son
argument appartient à la fonction et ne représente pas le produit d’une constante
indépendante par $\alpha$. Pour maintenir explicites les opérations qui
interviennent dans cette abréviation, on écrit
\begin{equation}
 B(\alpha)=\frac{9}{16}\alpha^2-\frac{5}{9}\alpha^3
 +\frac{7}{48}\alpha^4-\frac{1}{54}\alpha^5
 -\frac{90}{\pi}\mathcal C_\pi(\alpha),
 \label{md:compos:b-alpha}
\end{equation}
\begin{equation}
 \mathcal C_\pi(\alpha)=169\alpha^6+
 \frac{D_A\alpha^7}{1-D_A\alpha},
 \qquad D_A=\exp\!\left(-\frac{100\pi\alpha}{9}\right).
 \label{md:compos:correccion-regional}
\end{equation}
L’identité complète est
\begin{equation}
 C^*_{\deg}=\sqrt{\frac{1000\alpha}{\varphi}}+2B(\alpha).
 \label{md:compos:contraangulo-abreviado}
\end{equation}
En particulier, $\alpha$ appartient au radicand, et le terme radical est
conservé avec $2B(\alpha)$. Cette identité réécrit l’application construite
auparavant et maintient son objet et sa généalogie.

La définition originale conserve séparés le coefficient adimensionnel et la
section d’action :
\begin{equation}
 \eta_{\mathrm{ret}}=H_5-\frac{90}{\pi}\mathcal C_\pi,
 \qquad
 \hbar_{\mathrm{ret}}=\eta_{\mathrm{ret}}\,10^{-34}\mathcal U_S.
 \label{md:compos:seccion-accion}
\end{equation}
Par conséquent,
\begin{equation}
 \boxed{C^*_{\deg}=2\left(
 \frac{\hbar_{\mathrm{ret}}}{10^{-34}\mathcal U_S}+\alpha\right)
 =2(\eta_{\mathrm{ret}}+\alpha).}
 \label{md:compos:contraangulo-accion}
\end{equation}
L’expression utilise le coefficient adimensionnel de la section de Planck
réduite. La constante $h$ est obtenue ensuite par $h=2\pi\hbar$.
La composition consiste à extraire le coefficient d’action, à le composer avec
$\alpha$ et à appliquer la clôture bidirectionnelle. Ainsi est maintenu le type
dimensionnel de chaque opération.

La construction des masses exprime particule et antiparticule comme deux
sections orientées d’une même bande généalogique. L’opération
dodécaphasique est
\begin{equation}
 C_M(\tau)=-\operatorname{rev}(\tau),
 \qquad \tau_e=+++---+++---,
 \qquad C_M(\tau_e)=\tau_e,
 \label{md:compos:conjugacion-dodecafasica}
\end{equation}
prolongée au registre et à la représentation de charge. La route centrale est
fixe sous $C_M$, tandis que la représentation de charge distingue électron et
positron. Cette composition donne un contenu précis à la lecture bidirectionnelle.
Le facteur deux angulaire conserve sa fonction de clôture ; la conjugaison physique
conserve en outre les opérateurs indiqués. Le retour à $216$ restitue le
relèvement orienté : sa structure ne consiste pas à attribuer $108$
mises à jour à la particule et $108$ autres à l’antiparticule.

\subsubsection{Restitution radiale du coefficient d’action}
\label{md:compos:quinientos}

Nous écrivons $q_{\mathrm{rad}}=q_\varepsilon$ et $\epsilon=\varepsilon$ pour
le rapport radial et l’orientation de \eqref{rec:eq:eta-orientada}.
L’invariant radial orienté $\chi_{\mathrm{rad}}=\epsilon\log q_{\mathrm{rad}}$
permet d’écrire la restitution préalablement construite comme
\begin{equation}
 \frac{\eta_{\mathrm{ret}}}{\alpha}
 =-500\tanh\!\left(\frac{\chi_{\mathrm{rad}}}{2}\right)-1.
 \label{md:compos:restitucion-radial}
\end{equation}
Avec $\xi=C^*_{\deg}/A_{\deg}$ et $A_{\deg}=1000\alpha$, on obtient également
\begin{equation}
 \eta_{\mathrm{ret}}=\alpha(500\xi-1).
 \label{md:compos:restitucion-xi}
\end{equation}
Ce sont des opérateurs de récupération ultérieurs de la même section ; sa construction demeure
dans $H_5$ et dans le retour régional.

\subsection{Coordinatisations angulaires et conservation de l’action}
\label{md:compos:cartas}

La coordinatisation en degrés comporte un tour de $360$ unités et conserve les
partitions $12\cdot30=4\cdot90=3\cdot120=360$. La paire angulaire construite est
\begin{equation}
 A_{\deg}=1000\alpha,
 \qquad C^*_{\deg}=2(\eta_{\mathrm{ret}}+\alpha).
 \label{md:compos:par-angular}
\end{equation}
La conversion s’effectue une seule fois :
\begin{equation}
 x=\frac{\pi}{180}A_{\deg},\qquad
 y=\frac{\pi}{180}C^*_{\deg},\qquad
 q_\pm=\exp(-x\mp y).
 \label{md:compos:canales}
\end{equation}
Les exposants utilisent des représentants réels relevés et conservent le
nombre de tours. L’action satisfait
\begin{equation}
 \hbar\theta_{\mathrm{rad}}
 =\frac{h\theta_{\deg}}{360}.
 \label{md:compos:accion-cartas}
\end{equation}
Dans la fonctionnelle bilinéaire de Barbero, le transport de coordinatisation donne
\begin{equation}
 \frac{\pi A_{\deg}C^*_{\deg}}{180}
 =\frac{180xy}{\pi}.
 \label{md:compos:bilineal-cartas}
\end{equation}
La transformation bilinéaire fait partie de la même conversion qui est appliquée
aux arguments exponentiels ; l’omettre modifierait la fonctionnelle.

\subsection{Lecture électronique complète et retour d’action}
\label{md:compos:electron}

Les deux étapes de la même coordonnée électronique sont conservées :
\begin{equation}
 \varepsilon_e^{(0)}=\frac{\sqrt{3}}{4}
 \left[\exp\!\left(\frac{\varphi}{\pi^2}\right)-22\alpha^3\right]
 (1+15\Delta_4),\qquad
 \Delta_4=\frac{\pi-e\log\pi}{270}
 \left(1+\frac{\pi}{729}\right),
 \label{md:compos:electron-basal}
\end{equation}
\begin{equation}
 \boxed{\varepsilon_e^\beta=\varepsilon_e^{(0)}
 \left(\frac{H_5}{\eta_{\mathrm{ret}}}\right)^{80-54\alpha+6\Delta_4}.}
 \label{md:compos:electron-completo}
\end{equation}
Le centre $(5,5)$, son inversion, les déformations $(8,2)/(2,8)$ et le
défaut unitaire de quotient sont des antécédents de la valeur. La norme
$\sqrt{3}/4$ provient du commutateur entre projections additives et
multiplicatives ; le caractère exponentiel utilise l’orientation
$\varphi/\pi^2$ ; $22=27-5$ conserve un complément régional et $15=3\cdot5$
les incidences correspondantes. Les lecteurs de la route produisent
$(80,54,6)$. Le quotient $H_5/\eta_{\mathrm{ret}}$ incorpore le retour
d’action à la même section centrale.

La fonction du contre-angle est explicitée en substituant son inverse démontrée :
\begin{equation}
 \boxed{\varepsilon_e^\beta=\varepsilon_e^{(0)}
 \left(\frac{H_5}{C^*_{\deg}/2-\alpha}\right)^{80-54\alpha+6\Delta_4}.}
 \label{md:compos:electron-contraangulo}
\end{equation}
Cette expression conserve la génération électronique précédente. La même
section d’action qui participe au contre-angle détermine le facteur qui
fait passer la lecture basale à la lecture complète. Le quotient simplifie la base dimensionnelle
commune ; la réalisation absolue de masse conserve son lecteur temporel et sa coordinatisation
dimensionnelle. La substitution est une composition de dépendances construites,
sans introduire d’angle externe comme générateur électronique.

\subsection{Fonctionnelle de Barbero et transition hexade--octade}
\label{md:compos:barbero}

L’incidence maintient la même paire marquée lors de l’agrandissement du support :
\begin{equation}
 \mathcal F_6=\mathsf H\times\binom{\mathsf H}{2}
 \hookrightarrow
 \mathcal F_8=\mathsf O\times\binom{\mathsf H}{2},
 \qquad |\mathcal F_6|=90,\qquad |\mathcal F_8|=120.
 \label{md:compos:incidencia}
\end{equation}
Les cardinaux sont $6\cdot15$ et $8\cdot15$ ; le référent partagé permet de
comparer les deux incidences. La différence normalisée est $-1/12$. La chaîne
dodécaphasique conserve douze contributions hexadiques et une contribution
octadique de frontière de signe contraire. Le lecteur
\begin{equation}
 S_m(q)=\frac{q^m}{1-q^{3m}}
 =\sum_{j\geq0}q^{m+3mj}
 \label{md:compos:serie-retorno}
\end{equation}
conserve hauteur et retours. Ainsi se forment $12S_{90}-S_{120}$ et sa différence
entre les deux orientations. La fonctionnelle complète est
\begin{equation}
 \gamma=\frac{\pi A_{\deg}C^*_{\deg}}{180}
 +12\bigl[S_{90}(q_-)-S_{90}(q_+)\bigr]
 -\bigl[S_{120}(q_-)-S_{120}(q_+)\bigr].
 \label{md:compos:funcional-barbero}
\end{equation}
Barbero est, dans cette réalisation, l’évaluation orientée conjointe de la paire
angulaire et des retours de la transition d’incidences. Sa participation
à l’échelle d’aire se compose ensuite avec $\ell_P^2=G\hbar/c^3$ et avec
l’opérateur d’occupation. L’entropie d’intrication conserve en outre
l’état réduit et ses probabilités ; sa reconstruction nécessite ces données
et n’est pas déterminée par l’aire scalaire isolée.

\subsubsection{Élimination de la coordonnée commune électron--Barbero}
\label{md:compos:eliminacion}

Définissons $\kappa_e=80-54\alpha+6\Delta_4$. Sur la branche positive, avec
$\kappa_e\neq0$, l’équation électronique implique exactement
\begin{equation}
 \eta_{\mathrm{ret}}=H_5
 \left(\frac{\varepsilon_e^{(0)}}{\varepsilon_e^\beta}\right)^{1/\kappa_e}.
 \label{md:compos:eliminacion-eta}
\end{equation}
Pour exprimer la composition en maintenant visibles ses opérations, posons
$x=50\pi\alpha/9$, $y=\pi(\eta_{\mathrm{ret}}+\alpha)/90$ et
$\mathcal R(q)=12S_{90}(q)-S_{120}(q)$. Alors
\begin{equation}
 \gamma=\frac{100\pi}{9}\alpha(\eta_{\mathrm{ret}}+\alpha)
 +\mathcal R(e^{-x+y})-\mathcal R(e^{-x-y}).
 \label{md:compos:barbero-eta}
\end{equation}
La substitution de \eqref{md:compos:eliminacion-eta} dans cette identité produit
une relation explicite entre la lecture électronique complète et Barbero,
en conservant $\alpha$, $H_5$, $\Delta_4$ et les incidences. C’est une conséquence
de compatibilité de la construction. La branche physique angulaire utilisée
maintient $0<y<x$ ; l’opération n’introduit pas de masse métrologique comme donnée
génératrice.

Dans l’extension munie de l’étiquette d’orientation $\epsilon=\pm1$,
$C^*_{\deg,\epsilon}=2\epsilon(\eta_{\mathrm{ret}}+\alpha)$, l’échange
$\epsilon\mapsto-\epsilon$ laisse fixe $\eta_{\mathrm{ret}}$, conserve
l’expression électronique et permute $q_+$ et $q_-$. La fonctionnelle de Barbero change
de signe. Cette covariance distingue une lecture scalaire d’action d’une
évaluation orientée. L’identification à la conjugaison de charge conserve
en outre $C_M$ et la représentation électromagnétique des feuillets ; le signe angulaire
isolé ne remplace pas ces opérateurs.

\subsection{Moment orienté de Catalan et réponse constitutive}
\label{md:compos:catalan}

Le cycle dodécaphasique réalise un plan orienté de quart de tour. Son
coefficient résolvant est $k_4(s)=s/(1+s^2)$. Le moment de profondeur deux
est fixé par
\begin{equation}
 \mathcal C_2(s)=\sum_{k\geq0}
 \frac{(-1)^ks^{2k+1}}{(2k+1)^2},\qquad
 \mathcal D=s\frac{d}{ds},\qquad
 \mathcal D^2\mathcal C_2(s)=k_4(s),\qquad
 \mathcal C_2(1)=G_{\mathrm{Cat}}.
 \label{md:compos:momento-catalan}
\end{equation}
L’identité différentielle est utilisée pour $0<s<1$ ; la valeur extrême est
entendue avec les conditions de convergence et de restitution du moment déjà
établies. La définition intrinsèque de Catalan est la valeur extrême de ce
moment orienté de profondeur. Sa collection conserve plus d’information que
la valeur extrême. La réponse constitutive utilise cette collection :
\begin{equation}
 f(s)=\frac{1+s+s^2}{s(1+s)}\mathcal D^2\mathcal C_2(s)
 =\frac{1-s^3}{1-s^4},\qquad
 r_\pm=f(s_\pm),\qquad s_\pm=q_\pm^{30}.
 \label{md:compos:catalan-respuesta}
\end{equation}
Les puissances $3$ et $4$ correspondent à $90=3\cdot30$ et $120=4\cdot30$.
\begin{proposition}[Restitution fonctionnelle avec condition de frontière]
\label{md:compos:inversa-catalan}
La collection constitutive $f$ détermine le moment orienté par
\[
 \mathcal C_2(s)=\int_0^s\log\frac{s}{t}\,
 \frac{1+t}{1+t+t^2}f(t)\,dt,\qquad 0<s<1.
\]
C’est l’unique solution de
$\mathcal D^2\mathcal C_2(s)=s(1+s)f(s)/(1+s+s^2)$
qui tend vers zéro à l’origine.
\end{proposition}
\begin{proof}
Pour la collection construite,
$(1+t)f(t)/(1+t+t^2)=1/(1+t^2)$. L’intégrale converge, et sa
dérivée logarithmique est
$\mathcal D\mathcal C_2(s)=\int_0^s(1+t^2)^{-1}\,dt$.
Appliquer à nouveau $\mathcal D$ donne $s/(1+s^2)$. L’intégrale est
$O(s)$ à l’origine. Deux solutions diffèrent de $a+b\log s$ ;
la limite nulle impose $a=b=0$.
Pour $s<1$, intégrer la série géométrique sur des compacts et utiliser
$\int_0^s t^{2k}\log(s/t)\,dt=s^{2k+1}/(2k+1)^2$
retrouve la série de \eqref{md:compos:momento-catalan}.
Sa convergence absolue dominée permet $s\uparrow1$ et restitue Catalan.
\end{proof}
La condition de frontière fixe les deux constantes d’intégration.
La restitution utilise la fonction constitutive avec son domaine ; la
valeur extrême de Catalan et les deux valeurs de canal sont des lectures de
cette fonction déjà construite. À l’extrémité, l’identité doublement
dérivée conserve son interprétation par limite d’Abel.

L’opérateur constitutif $V=r_+P_++r_-P_-$ détermine
\begin{equation}
 \widehat\varepsilon=r_+^2,\qquad \widehat\mu=r_-^2,
 \qquad \widehat Z=\frac{r_-}{r_+},\qquad
 \widehat c=(r_+r_-)^{-1}.
 \label{md:compos:lectores-constitutivos}
\end{equation}
Avec $\psi=f^{-1}$ et
\begin{equation}
 \mathcal H_B(s)=\frac{12s^3}{1-s^9}
 -\frac{s^4}{1-s^{12}}-\frac{(\log s)^2}{20\pi},
 \label{md:compos:h-barbero}
\end{equation}
la factorisation construite prend la forme
\begin{equation}
 \gamma=\mathcal H_B(\psi(r_-))-\mathcal H_B(\psi(r_+)).
 \label{md:compos:factorizacion-barbero}
\end{equation}
La partie logarithmique restitue exactement la contribution angulaire de
Barbero, conservée en degrés ou transportée intégralement en radians.

\subsection{Restitution par la vitesse normalisée et Barbero}
\label{md:compos:difeomorfismo}

\paragraph{Proposition.}
Dans la collection constitutive normalisée à deux canaux étiquetés, l’application
\begin{equation}
 (r_+,r_-)\longmapsto(\widehat c,\gamma),\qquad
 (3/4,1)^2\longrightarrow(1,16/9)\times\mathbb R
 \label{md:compos:mapa-global}
\end{equation}
est un difféomorphisme. Son inverse récupère les deux valeurs propres étiquetées.
Les projecteurs de feuillet appartiennent à la représentation de départ : deux
scalaires ne récupèrent pas une base spatiale arbitraire.

\paragraph{Démonstration.}
L’expression régulière
$f(s)=(1+s+s^2)/(1+s+s^2+s^3)$ donne
\begin{equation}
 f'(s)=-\frac{s^2(s^2+2s+3)}{(1+s+s^2+s^3)^2}<0.
 \label{md:compos:derivada-f}
\end{equation}
Ainsi, $f$ est analytique et strictement décroissante, avec $f(0^+)=1$ et
$f(1^-)=3/4$. Pour $0<s<1$,
\begin{equation}
 \mathcal H_B'(s)=
 \frac{36s^2(1+2s^9)}{(1-s^9)^2}
 -\frac{4s^3(1+2s^{12})}{(1-s^{12})^2}
 -\frac{\log s}{10\pi s}>0.
 \label{md:compos:derivada-h}
\end{equation}
En effet, le quotient du premier terme positif par le second terme
positif avant de le soustraire est
\begin{equation}
 \frac{9}{s}\frac{1+2s^9}{1+2s^{12}}
 \left(\frac{1-s^{12}}{1-s^9}\right)^2>9.
 \label{md:compos:cota-derivada-h}
\end{equation}
Le dernier terme de \eqref{md:compos:derivada-h} est également positif.
Par conséquent, $F_B=\mathcal H_B\circ\psi$ possède une dérivée strictement négative,
$F_B(3/4^+)=+\infty$ et $F_B(1^-)=-\infty$.

$c=\widehat c$ étant fixé, nous écrivons $z=\log\widehat Z$ et
$r_\pm=c^{-1/2}\exp(\mp z/2)$. L’intervalle admissible est
\begin{equation}
 |z|<a(c):=\min\left\{\log c,\log\frac{16}{9c}\right\}.
 \label{md:compos:dominio-z}
\end{equation}
La fonction $\Gamma_c(z)=F_B(r_-)-F_B(r_+)$ satisfait
\begin{equation}
 \Gamma_c'(z)=\frac12
 \bigl[r_-F_B'(r_-)+r_+F_B'(r_+)\bigr]<0.
 \label{md:compos:derivada-gamma}
\end{equation}
Lorsque $z$ augmente jusqu’à $a(c)$, soit $r_-$ atteint $1$, soit $r_+$ atteint $3/4$,
soit les deux se produisent ; $\Gamma_c$ tend vers $-\infty$. À l’extrémité opposée,
elle tend vers $+\infty$. Il existe donc une unique solution pour chaque $\gamma$
réel. La dérivée non nulle et le théorème des fonctions implicites fournissent
une inverse lisse globale. Cela prouve l’énoncé. Dans l’orientation
originale $y>0$, on a $z<0$ et $\gamma>0$.

La vitesse normalisée conserve le produit des réponses ; Barbero
distingue de manière monotone leur répartition orientée. Ensemble,
ils restituent l’information spectrale que chaque lecture isolée ne détermine pas.

L’énoncé concerne la collection des lecteurs constitutifs. L’image
effectivement produite par le générateur HMT complet est un sous-ensemble de
cette collection : la restriction à cette image conserve l’injectivité et l’ordre
causal. La normalisation et les marques de feuillet sont des données nécessaires de la
restitution. La preuve n’introduit pas de valeurs CODATA pour choisir les canaux ni
n’affirme que chaque point synthétique de la collection soit un état engendré.

\subsection{Restitution sur l’image engendrée}
\label{md:compos:restitucion-completa}

Une fois $r_\pm$ récupérés à partir de $(\widehat c,\gamma)$, on obtient
$s_\pm=\psi(r_\pm)$. Alors
\begin{equation}
 A_{\deg}=-\frac{3}{\pi}\log(s_+s_-),\qquad
 C^*_{\deg}=\frac{3}{\pi}\log\frac{s_-}{s_+}.
 \label{md:compos:restitucion-angular}
\end{equation}
Sur l’image engendrée et dans son orientation positive,
\begin{equation}
 \alpha=\frac{A_{\deg}}{1000},\qquad
 \eta_{\mathrm{ret}}=\frac{C^*_{\deg}}{2}-\alpha>0.
 \label{md:compos:restitucion-alpha-eta}
\end{equation}
Avec les coordonnées régionales $\pi,\varphi,e$ et les sections dimensionnelles
déjà construites, cette restitution redonne $H_5$, le facteur
$R_{\mathrm{act}}=H_5/\eta_{\mathrm{ret}}$ et la lecture électronique complète
de la section~\ref{md:compos:electron}. Elle redonne également la section de
Planck. Si l’on conserve en outre la base positionnelle à douze blocs, l’origine
et la retenue, l’identité
\begin{equation}
 \kappa_{\mathrm{ag}}=\pi+e-\varphi-4-\alpha
 \label{md:compos:restitucion-k}
\end{equation}
permet la restitution du registre $K$ décrite par son lecteur. Il s’agit
d’un parcours de reconstruction ; son sens inverse ne constitue pas une nouvelle
flèche génératrice à partir de valeurs physiques externes.

La gravité et la conversion thermique conservent les sections du même
état. Le théorème~\ref{md:rigidez:teorema} et
\eqref{md:rigidez:reloj} déterminent d’abord $L_\alpha$ et sa
représentation circulaire $R_{108}=L_\alpha$ :
\begin{equation}
 G_a=\frac{c_{\mathrm{int}}^3R_{108}^2}{\hbar_a},\qquad
 \frac{G_a\hbar_a}{c_{\mathrm{int}}^3}=R_{108}^2,
 \qquad
 k_B\Theta_{\mathrm{clk},a}\log3=\frac{h_a}{108t_0}.
 \label{md:compos:gravedad-termica}
\end{equation}
Le rayon chronologique, la vitesse physique et la base thermique conservent leurs
constructeurs. En particulier, l’horloge de la représentation circulaire est
liée à $L_\alpha$ et $c_{\mathrm{int}}$. La première
composition explicite la réciprocité action--gravité ; la seconde réunit
action, durée et information tritique, avec individualisation de la température
et de $k_B$ dans leurs bases déclarées.

\subsection{Définitions structurelles des lectures composées}
\label{md:compos:definiciones}

\begin{description}
 \item[$\alpha$.] Coordonnée scalaire de clôture dodécaphasique compatible avec
 les régions et avec le registre arithmétique--géométrique $K$ ; elle transmet cette
 compatibilité au secteur angulaire et d’action.
 \item[Contre-angle.] Coordonnée angulaire de la clôture bidirectionnelle qui
 combine le coefficient de la section d’action de retour avec $\alpha$,
 dans la coordinatisation de $360$ unités.
 \item[Planck réduite.] Section d’action dont le coefficient régional et l’échelle
 déterminantale sont construits avant le contre-angle ; dans l’ellipse, mesure
 d’aire orientée par unité de phase paramétrique.
 \item[Électron.] Section centrale persistante ; sa lecture scalaire compose
 incompatibilité additive--multiplicative, transfert régional, défaut
 de clôture et retour d’action. L’équation complète conserve la section,
 en plus de sa valeur de masse.
 \item[Barbero--Immirzi.] Fonctionnelle orientée du transport angulaire et des
 retours de l’incidence hexade--octade ; dans la collection constitutive,
 coordonnée qui, avec la vitesse normalisée, restitue le spectre
 étiqueté.
 \item[Catalan.] Valeur extrême du moment orienté de profondeur deux du
 quart de tour dodécaphasique. Sa collection est reliée différentiellement à la
 réponse du vide.
 \item[$G$.] Coefficient commun de la réalisation radiale qui relie
 énergie et longueur conjuguée en conservant simultanément
 $H\Lambda=\hbar cI$ et $R\Lambda=L_\alpha^2I$.
 La composition longitudinale et polaire détermine sa normalisation ;
 son expression dimensionnelle est $c^3L_\alpha^2/\hbar$.
 \item[Boltzmann.] Transducteur entre la droite thermique et l’énergie de
 décision, avec une normalisation informationnelle tritique et une base thermique déclarées.
\end{description}
Ces formulations intrinsèques se fondent sur les objets et équations
construits. Leur portée est celle de cette réalisation ; les identifications
physiques conservent leurs démonstrations correspondantes et ne se déduisent pas de
la seule dénomination structurelle.

\subsection{Transport du bloc électronique nilpotent}
\label{md:compos:nilpotentes}

\paragraph{Proposition.}
Soient
\begin{equation}
 S=\begin{pmatrix}1&0\\0&-1\end{pmatrix},\qquad
 J=\begin{pmatrix}0&1\\-1&0\end{pmatrix},\qquad
 n_\pm=\frac{S\pm J}{2},\qquad
 D=\operatorname{diag}(r,r^{-1}),\quad r>0.
 \label{md:compos:bloque-nilpotente}
\end{equation}
Définissons $n_{\pm,D}=Dn_\pm D^{-1}$ et $g=D^{-T}D^{-1}$. Alors
\begin{equation}
 n_{\pm,D}^2=0,\qquad
 \{n_{+,D},n_{-,D}\}=I,\qquad
 (n_{+,D})^{\dagger_g}=n_{-,D},
 \label{md:compos:identidades-nilpotentes}
\end{equation}
où $A^{\dagger_g}=g^{-1}A^Tg$. L’ellipse transportée depuis
$Y^TY=2\hbar$ au moyen de $X=DY$ satisfait $X^TgX=2\hbar$.

\paragraph{Démonstration.}
Le calcul matriciel direct donne
$S^2=I$, $J^2=-I$ et $SJ+JS=0$. Ainsi,
\[
 n_\pm^2=\tfrac14(S^2+J^2\pm SJ\pm JS)=0,
 \qquad
 n_+n_-+n_-n_+=\tfrac12(S^2-J^2)=I.
\]
En outre, $S^T=S$ et $J^T=-J$ impliquent $n_+^T=n_-$. La conjugaison par $D$
préserve produits, carrés et identité. Pour l’adjoint métrique,
\[
 g^{-1}(Dn_+D^{-1})^Tg
 =DD^T D^{-T}n_-D^T D^{-T}D^{-1}
 =Dn_-D^{-1}=n_{-,D}.
\]
Enfin, $D^TgD=I$, d’où
$X^TgX=Y^TD^TgDY=Y^TY=2\hbar$. Les affirmations sont ainsi prouvées.

La vérification algébrique exacte a lieu dans
$\mathbb Q[r,r^{-1}]$ ; la restriction $r>0$ fixe la réalisation elliptique
orientée utilisée. Cette composition conserve les fonctions particulières
des différents facteurs $1/2$ du développement : une coïncidence numérique
entre eux ne constitue pas une identification de leurs opérateurs.
